\begin{table}[h]
\centering
\caption{Main Results: Kaitz-Weighted Minimum Wage Effect on Log Employment}
\label{tab:main}
\tiny
\begin{tabular}{lccccc}
\hline\hline
& (1) & (2) & (3) & (4) \\
& Baseline & +Controls & Alternative & Sample \\
& & & Kaitz (2016) & Split \\
\hline
Kaitz Index × MW Change & $-0.084^{**}$ & $-0.081^{**}$ & $-0.079^{*}$ & $-0.083^{*}$ \\
& (0.032) & (0.034) & (0.041) & (0.044) \\
\hline
Sector FE & Yes & Yes & Yes & Yes \\
Quarter FE & Yes & Yes & Yes & Yes \\
Pre-2012 employment level & No & Yes & Yes & No \\
Sector employment growth 2008--2012 & No & Yes & Yes & No \\
\hline
Observations & 1408 & 1408 & 1408 & 704 \\
R-squared & 0.834 & 0.847 & 0.835 & 0.841 \\
Clusters (sectors) & 22 & 22 & 22 & 11 \\
\hline
\multicolumn{5}{l}{\textit{Implied Elasticities:}} \\
\quad Sector with Kaitz = 0.50 & -0.042 & -0.041 & -0.040 & -0.042 \\
\quad Sector with Kaitz = 0.20 & -0.017 & -0.016 & -0.016 & -0.017 \\
\hline\hline
\end{tabular}
\begin{flushleft}
\small \textit{Notes:} Each cell reports the coefficient on the interaction between pre-2012 Kaitz index and log cumulative minimum wage increase from 2012 baseline. Standard errors are clustered by sector (22 clusters). *** $p<0.01$; ** $p<0.05$; * $p<0.10$. Column 4 restricts sample to 2012--2020 (pre-formula period) for robustness. Negative coefficients indicate that high-exposure sectors (high Kaitz) experience larger employment declines as the minimum wage rises, controlling for sector and time fixed effects.
\end{flushleft}
\end{table}
